\section{Base Change}
We have stated before that the information about a ring $R$ is encoded in a deep way into $\Rmod{R} $, and now we can generalize this notion using the language of functors to also study homomorphisms $R\to S$, by rephrasing this notion into that of functors $\Rmod{R} \to \Rmod{S} $.
\subsection{Balanced Maps}
Tensor products satisfy an even stronger universal property than previously explored.
\begin{definition}[Balanced]\label{dfn:8.11}
	Let $M,N$ be $R$-modules for $R$ commutative, and let $G$ be a $\mathbb{Z} $-module; that is, $G$ is an abelian group. A $\mathbb{Z} $-bilinear map $\phi : M\times N \to G$ is $R$-\emph{balanced} if for all $m\in M$, $n\in N$, and $r\in R$,
	\[
		\phi (r m,n)=\phi (m,rn)
	.\]
\end{definition}

Defining such a notion with $G$ an $R$-module would make all bilinear maps $R$-balanced, as $\phi (r m,n)=r\phi (m,n)=\phi (m,rn)$ due to bilinearity. However, since $G$ is not an $R$-module, this is a more general case. It turns out $R$-balanced maps factor uniquely through the tensor product as well.
\begin{lemma}\label{lma:8.3}
	Let $R$ be a commutative ring, let $M,N$ be $R$-modules, and let $G$ be an abelian group. Then every $\mathbb{Z} $-bilinear, $R$-balanced map $\phi : M\times N \to G$ factors uniquely through $M\otimes_R N$; that is, there exists a unique group homomorphism $\overline{\phi } : M\otimes_R N \to G$ such that the diagram
	\[\begin{tikzcd}[row sep = 2.5em, column sep = 2.5em]
		M\times N\arrow[r, "\phi "]\arrow[d, "\otimes "']&G\\
		M\otimes _RN\arrow[ru, dashed, " \overline{\phi } "']
	\end{tikzcd}\]
	commutes.
\end{lemma}
\begin{proof}
	Recall that we constructed $M\otimes _RN$ precisely as a quotient
	\[
		\frac{R^{\oplus (M\times N)} }{K} 
	,\]
	where $K$ is the $R$-module generated by the relations necessary to make the map $M\times N\to R^{\oplus (M\times N)} /K$ $R$-bilinear. That is,
	\[
		M\times N\to R^{\oplus (M\times N)} \to \frac{R^{\oplus (M\times N)} }{K} 
	.\]
	We also saw that every element of $M\otimes _RN$ can be written as a linear combination of pure tensors
	\[
		\sum_{i} m_i\otimes n_i
	.\]
	We can then define a surjective map $\mathbb{Z} ^{\oplus (M\times N)} \to M\otimes _RN$ by $(m,n)\mapsto m\otimes n$. The kernel $K_B$ of this group homomorphism is all elements that map to $0$ in $R$-bilinear maps; that is, generated by elements of the form
	\[
		(m,n_1+n_2)-(m,n_1)-(m,n_2)
	,\]
	\[
		(r m,n)-(m,rn)
	,\]
	and similarly in $m$ for the first one. We then have an isomorphism of abelian groups
	\[
		\frac{\mathbb{Z} ^{\oplus (M\times N)} }{K_B} \cong \frac{R^{\oplus (M\times N)} }{K} 
	.\]
	By the same logic as the tensor product, we have that $M\otimes_R N$ uniquely factors $\mathbb{Z}$-bilinear $R$-balanced maps.
\end{proof}

Taking $\phi $ to be $R$-linear recovers the universal property discussed in the previous section. Balanced maps can in fact be defined if $M$ is a left $R$-module and $N$ is a right $R$-module, but the issue is that $M\otimes _RN$ does not necessarily carry $R$-module structure in this case.

\subsection{Bimodules; Adunction Again}
This new universal property allows us to upgrade the adjunction of tensor, but this requires defining bimodules.
\begin{definition}[Bimodule]\label{dfn:8.12}
	Let $R$ and $S$ be commutative rings. An $(R,S)$-bimodule is an abelian group $N$ with compatabile $R$- and $S$-module structures; that is, for all $n\in N$, $r\in R$, and $s\in S$,
	\[
		r(sn)=s(rn)
	.\]
	One could also say $N$ is a left $R$-module and a right $S$-module, and give $(rn)s=r(ns)$.
\end{definition}

For example, every $R$-module is an $(R,R)$-bimodule for $R$ commutative. If $M$ is an $R$-module and $N$ is an $(R,S)$-bimodule, then $M\otimes _RN$ obtains an $S$-module structure defined on pure tensors by
\[
	s(m\otimes n)=m\otimes (sn)
.\]
We thus obtain an $(R,S)$-bimodule structure on $M\otimes _RN$.
